\section{Related Work}\label{s:related-work}
% 3 pages


\subsection*{Transformer Architecture and Representations}

We use the standard transformer architecture \citep{vaswani2017attention} 
with attention-only layers (no MLPs). Each layer computes a weighted sum of value vectors, where the weights are given by a softmax over query-key dot products. 
Throughout this paper we use the framework from \citet{elhage2021mathematical} to discuss transformers, and refer the reader there and to \citet{ferrando2024primer} for more details. Figure \ref{fig:rq2_pipeline} provides a visual overview. This depicts the residual stream as a shared vector across all layers which read from and write to it across depth. The output is the final residual stream state. This framing enables us to easily attribute specific information to specific layers and positions.

We use the term ``feature'' to mean a property of an input that a neural network may learn to represent internally, as in \cite{wattenberg2024relational}. However, we note that there is no consensus definition in the literature \cite{elhage2022toy,wattenberg2024relational}.
Our model takes ordered symbol pairs (bigrams) as input, and we treat the individual symbols and their ordering as three separate features. Since neural networks have a black-box nature, we cannot directly observe which features a model has learned. Mechanistic interpretability addresses this by attempting to reverse-engineer those representations from model parameters and behaviour.

Neural networks often contain ``polysemantic'' neurons, which respond to multiple seemingly unrelated features (``monosemantic'' neurons correspond to only one) \citep{olah2020zoom}. One account of this is ``superposition'' \citep{elhage2022toy}: because 
real-world features tend to be sparse (rarely co-occurring), a model 
can represent more features than its activation dimension allows by 
assigning each feature a nearly orthogonal direction in the residual 
stream, accepting occasional interference when rare co-activations 
occur. \citet{elhage2022toy} demonstrate this with a toy model that 
encodes five sparse features in two dimensions. This motivates 
the use of SAEs: if polysemantic activations arise from superposed 
sparse features, then a dictionary learning method that enforces 
sparsity should be able to recover those features as separate latents. We refer the reader to \cite{faruqui-etal-2015-sparse, arora-etal-2018-linear, Subramanian_Pruthi_Jhamtani_Berg-Kirkpatrick_Hovy_2018, yun-etal-2021-transformer} for more information on dictionary learning outside of interpretability.

% In our paper, we use the word ``feature'' to mean a property of an input to a neural network. Our transformer model takes bigrams as input, and we refer to each element and their ordering in the bigram as separate features.
% To generalise to a given dataset (rather than memorise it), models must learn to represent some or all of these features in their latent space. However, since they have a black-box nature, we do not know which features they have learned to represent during training. 
% Hence, we can aim to reverse engineer the model from its parameters and behaviours to interpret what features the model has learned.
% Individual neurons typically do not represent individual features (monosemanticity), which would aid reverse-engineering because we could map each neuron to a specific role in the computation. Instead, they exhibit polysemanticity, which means a single neuron represents multiple features. One hypothesised reason for this is superposition \cite{elhage2022toy}: the idea that models can represent more features than their dimensionality allows by representing more than one feature in a single linear dimension. For example, \citet{elhage2022toy} show that a toy model with five sparse (i.e. rarely co-appearing) input features could represent them using just two linear dimensions by mapping each feature to five nearly orthogonal directions in 2D space.
% Intuitively, it's plausible that models trained on real-world data exhibit superposition, because most real-world features are sparse. In that case, we could encode a single polysemantic neuron's activations, which are feature-dense, using (possibly many) more monosemantic neurons (which each represent a single feature).
% SAEs have been proposed to attempt just that.

\subsection*{Mechanistic Interpretability}
Mechanistic interpretability is the study of reverse-engineering neural networks 
into human-interpretable algorithms \citep{olah2020zoom, 
ferrando2024primer, nanda2023progress}. This is usually achieved by identifying ``circuits'' in the model: paths of computation through a model that are responsible for specific behaviours, represented as model subgraphs.
Some example behaviours that have been successfully reverse-engineered in this way are indirect object identification \citep{wang2022interpretability}, in-context learning via induction heads \citep{olsson2022context}, and modular arithmetic via Fourier representations \citep{nanda2023progress}. 
Many of these studies use toy models as a controlled environment for identifying mechanisms that might generalise to larger models, which supports our use of toy models in this paper. For example, \citet{baeumel2025modular} found the modular arithmetic results on toy models subsequently generalised to pretrained LLMs.

Moreover, recent work has found ``steering'' vectors in the residual stream which can be steered via linear scaling and patched back into the residual vector to predictably change model behaviour. For example,
\citet{arditi2024refusal} use the difference-in-means vector between residual streams on harmful and harmless instructions to identify a single ``refusal direction'' in the residual stream of safety-tuned LMs, and use it to to suppress or amplify refusal behaviour via linear scaling.
% adding negative multiples of a direction correlated with a given output sentiment ($\boldsymbol u$) to the residual stream ($\boldsymbol r{L\_i}' \leftarrow \boldsymbol r{L\_i} - \alpha \boldsymbol u$) can cause the model to generate output with the opposite sentiment \cite{tigges2024language}
This provides evidence that a given direction in the residual stream is causally relevant for a given behaviour, which provides insights into what feature that direction represents \cite{ferrando2024primer}.
We use the same principle in our latent steering experiments, in which we steer SAE latents rather than raw residual stream directions and patch the SAE-reconstructed activations into the residual stream.

% When most technologies are built, their capabilities are not a surprise because they were specifically implemented. Neural networks (including transformers) are not like this because their capabilities emerge from various training, fine-tuning and reinforcement learning processes. Therefore, neural networks tend to have a ``black-box'' nature, where its creators do not necessarily understand the internal processes that have emerged from gradient descent.

% Mechanistic interpretability aims to remove this black-box nature by reverse-engineering neural networks into human-interpretable algorithms \cite{olah2020zoom,ferrando2024primer,nanda2023progress}. 
% Recent work has demonstrated success across several areas. One major direction involves localizing specific model behaviours to discrete subgraphs, or "circuits". For example, researchers have reverse-engineered the circuits responsible for indirect object identification (IOI) \cite{wang2022interpretability}, the "greater-than" logical operation \cite{ashioya2025greaterthan}, and "induction heads," which drive in-context learning \cite{olsson2022context}.

% Beyond isolating circuits for specific tasks, mechanistic interpretability has explained broader structural phenomena within neural networks. Investigations into superposition have shown how models represent more features than their activation dimensions allow by compressing features into nearly orthogonal directions \cite{elhage2022toy}. Additionally, research has identified linear emergent world models, demonstrating that models trained on sequence prediction (such as Othello moves) develop linear representations of the underlying state space \cite{li2023emergent, nanda2023othello}.

% The field has also provided explanations for training dynamics. By tracking circuit formation over time, researchers have explained phenomena such as "grokking" in algorithmic tasks, where models transition from memorization to generalization by learning understandable functions like Fourier transforms for modular addition \cite{nanda2023progress}.

% Furthermore, mechanistic interpretability supports direct intervention in model behaviour. By isolating specific activation vectors, such as a single "refusal direction" in chat-tuned models, researchers can steer outputs and bypass safety filters without additional fine-tuning \cite{arditi2024refusal}.

\subsection*{Relational composition and binding}

\begin{table*}[t]
\centering
\caption{Comparison of relational composition mechanisms and their implications for interpretability.}
\label{tab:composition_mechanisms}
\begin{tabular}{p{3.5cm}p{4cm}p{4cm}p{4cm}}
\toprule
\textbf{Composition} & \textbf{Description} & \textbf{Example} & \textbf{Implications} \\
\midrule
Direction-based &
Features are vectors; sequence position defined by different directions.
& Matrix binding \cite{wattenberg2024relational} - features projected to distinct linear subspaces: $r = Ax + By$. &
Feature multiplicity: SAEs may learn echo features ($Ax$, $Ay$, $Bx$, $By$) representing the same underlying concept. \\
\addlinespace
Fixed magnitude-based &
Multiple features have same direction; sequence position defined by different magnitudes.
&
Onion-like representations \cite{csordas-etal-2024-recurrent} - ``unpeel'' via autoregression to retrieve sequence positions.
&
Counter-example to the widely held assumption that representations are linear.
\\
\addlinespace
Relative magnitude-based &
Sequence position defined by magnitude relative to other positions in sequence.
& This work - sequence order defined by a weighted sum of input token and positional embeddings. &
Counter-example to the widely held assumptions that representations are linear and SAE latents are binary.
\\
\bottomrule
\end{tabular}
\end{table*}

Classic work in distributed representations show how single vectors can encode structured data. Tensor Product Representations (TRPs) and Vector Symbolic Architectures (VSAs) bind and additively superpose vectors, enabling ordered structures in fixed width vectors \citep{smolensky1990tpr, plate1995hrr, kanerva2009hdc}.

\citet{wattenberg2024relational} highlight additive matrix binding, where vectors are written to slots using role-specific linear maps so that bigrams $(x, y),(y,x)$ remain separable, as a candidate mechanism for relational composition in transformers. That is, given an ordered bigram $(d_1, d_2)$ and two $n \times n$ matrices $A$ and $B$, define the representation of $(x, y)$ by
\(
r = Ax + By
\).
In this way, $(x, y)$ and $(y, x)$ have different representations.
\citet{wattenberg2024relational} claim that this causes \textit{feature multiplicity}: a kind of false positive where multiple ``echo features'' are created that represent the same concept in different contexts. Given two features $x$ and $y$, matrix binding results in $Ax$, $Ay$, $Bx$, and $By$ features that represent the same concept as $x$ and $y$. 

Empirically, \citet{feng2023language, prakash2023fine} identify binding ID vectors that bind entities to their attributes in-context in language models. For example, a ``green car'' is represented by a vector representing ``green'', a vector for ``car'', and a binding ID vector to join the two. A binding ID vector $k$ is attached to an entity and its attribute by addition in their respective activation spaces. To answer a query for a given entity, the attribute with a matching binding ID is retrieved from the context activation space. \citet{feng2023language} also show that these mechanisms are position-independent with respect to the attribute and entity, and that the activations can be factorised.

\citet{brinkmann2024mechanistic} reverse-engineer an attention-only tree-search transformer that stores precomputed subpaths in special token positions and merges them later. This occurs when the tree depth, $N$, is greater than $L-1$, where $L$ is the number of transformer layers. This constraint is also necessary in our work, where $N$ is the size of the input $N$-gram. However, the authors do not investigate the structure of these activations.

\citet{csordas-etal-2024-recurrent} demonstrate ``onion representations'' in small gated recurrent neural networks (RNNs), where these slots are defined by different orders of magnitude rather than distinct linear subspaces.
The model represents ordered sequences by closing the gates gradually and synchronously over the input phase, which exponentially decays the scaling factor of subsequent sequence positions.
This produces layered onion representations where earlier sequence positions dominate the magnitude space.
This contrasts larger models which indicate sequence position by sharply closing their gates, which creates position-dependent subspaces for each input.
The onion representation allows autoregressive decoding to sequentially ``peel'' off tokens by subtracting the current dominating token embedding.
Multiple tokens can occupy the same directional subspace at different magnitude scales; any linear direction will cross-cut multiple layers of the onion. 
While the authors find that sequence order is represented by fixed magnitude scales, we find an alternative mechanism where it is instead represented by relative magnitude between positions, which is input-dependent and not fixed.

Our toy model provides evidence for magnitude-based composition in attention-only transformers, but with \textit{relative} magnitudes rather than fixed ones as observed in \citet{csordas-etal-2024-recurrent}. 
% Rather than each position being assigned a predetermined magnitude scale $\gamma^t$, the compositional information is encoded in the relative scaling between feature components that can adapt to the specific features being combined. Additionally, unlike GRUs which require autoregressive unpeeling to decode onion representations, transformers can potentially access multiple positions in parallel through attention mechanisms.
See Table \ref{tab:composition_mechanisms} for an overview of these composition methods.

\subsection*{Sparse Autoencoders} 
Models can store more sparse features than the dimension of their activations naively allows through superposition, creating challenges for interpreting these models \citep{elhage2022toy}. Sparse Autoencoders (SAEs) \citep{bricken2023monosemanticity, huben2024sparse} and their variants \citep{gao2024scaling, bussmann2024batchtopk, rajamanoharan2024jumping, leask2025inference, costa2025evaluating, bussmann2025learning} have been proposed as a tool for recovering these sparse features from dense activations. 
An SAE is typically implemented as an autoencoder with a single hidden layer, trained to reconstruct model activations while keeping the hidden representation sparse using a penalty (such as an $L_1$ regularization term) \citep{huben2024sparse,bricken2023monosemanticity}. 
The encoder maps dense activation to a high-dimensional sparse latent space, and the decoder reconstructs the original activations using a learned dictionary of monosemantic latents, which ideally corresponds to interpretable features.
% The motivation is that if polysemantic activations arise from superposition, an SAE should decompose them into monosemantic latents, each corresponding to one feature.
% The encoder projects the dense activations into a high-dimensional sparse latent space, and the decoder reconstructs the original activations using a dictionary of feature vectors. 
% An SAE is an autoencoder with a single hidden layer that is trained to reconstruct its input while enforcing sparsity in the activations of the neurons in this hidden layer \cite{huben2024sparse,bricken2023monosemanticity}.


\begin{table*}[t]
\centering
\caption{Comparison of some sparse autoencoder (SAE) variants.}
\label{tab:sae_types}
\begin{tabular}{p{2.5cm}p{12.5cm}}
\toprule
\textbf{Type} & \textbf{Description} \\
\midrule
Standard ReLU \cite{bricken2023monosemanticity}
  & Applies an $L_1$ penalty to hidden activations using a standard ReLU activation function to encourage sparsity. \\
\addlinespace
Gated \cite{rajamanoharan2024improving}
  & Decouples feature detection from magnitude estimation using a gated mechanism to prevent the shrinkage issues caused by standard $L_1$ penalties. \\
\addlinespace
JumpReLU \cite{rajamanoharan2024jumping}
  & Uses a discontinuous step function at a learned per-feature threshold, approximating an $L_0$ penalty through straight-through estimation. Notably used in \cite{lieberum2024gemma}. \\
\addlinespace
BatchTopK \cite{bussmann2024batchtopk}
  & Retains top $n \times k$ activations over a batch of $n$ inputs and sets the rest to zero, where $k$ is a tuned hyperparameter. This forces an average of $k$ latent activations per input, so the actual $L_0$ sparsity is approximately $k$. Notably used in \cite{Brixi2026, karvonen2025saebench}. \\
\addlinespace
Matryoshka \cite{bussmann2025learning}
  & Trains nested sub-dictionaries of increasing size that share a single encoder/decoder, producing multiple resolutions of feature granularity in one run. Notably used in \cite{mcdougall2025gemmascope2}. \\
\bottomrule
\end{tabular}
\end{table*}

Several architectural variants address limitations in the standard ReLU-based SAE (Table \ref{tab:sae_types}). For example, standard $L_1$ penalties often cause "shrinkage," where the reconstructed feature magnitudes are systematically underestimated \cite{ferrando2024primer}. Gated SAEs \cite{rajamanoharan2024improving} address this by decoupling feature detection from magnitude estimation using a gated mechanism. JumpReLU SAEs \citep{rajamanoharan2024jumping} replace the standard ReLU with a discontinuous step function at a learned threshold, optimizing an $L_0$ penalty to achieve the same decoupling.
TopK SAEs \citep{gao2024scaling} replace the sparsity penalty by explicitly retaining only the $k$ highest pre-activations per token. BatchTopK SAEs \citep{bussmann2024batchtopk} extend this by enforcing the TopK constraint across an entire batch, allowing natural variation in the number of active latents per token while maintaining the target sparsity in expectation. Matryoshka SAEs \citep{bussmann2025learning} train nested sub-dictionaries of increasing size under a shared encoder and decoder, producing representations at multiple levels of granularity from a single training run. We use Matryoshka, BatchTopK and JumpReLU variants in this paper.

% explain SAEs - https://arxiv.org/pdf/2412.06410

SAE performance is typically evaluated as the pareto frontier of two metrics: $L_0$ norm of the SAE latent activations vector, which measures sparsity, and reconstruction quality (explained variance, cross-entropy increase when patching reconstructions into the model) \citep{karvonen2025saebench,ferrando2024primer,bricken2023monosemanticity}.

SAEs have been successfully used on exploratory interpretability tasks like model auditing \citep{marks2025auditing} and hypothesis generation \citep{movva2025sparse}, but it is unclear whether they can be used to recover the true features of models \citep{leask2025sparse, chanin2024absorption}.
For example, \citet{leask2025sparse} show that SAEs can achieve high reconstruction quality while learning latents that do not correspond to any interpretable unit of computation.
Therefore, reconstruction metrics are necessary but not sufficient evidence that an SAE has recovered the model's true representational structure.

\citet{wattenberg2024relational} argue that matrix binding will result in echo features in SAEs, where the SAE learns latents corresponding to projections of the same feature into different binding subspaces. The effects of ID vectors and onion representations on SAEs have so far not been studied. 

SAE latents are generally treated as binary features in the literature \citep{paulo2024automatically, quirke2025binary}, where they are 'on' if the feature is active and 'off' otherwise. In contrast, our results demonstrate graded latent activations in which relative activation magnitude encodes relational information, such as sequence order. If these results generalise to LLMs, then they could undermine this binary perspective. Additionally, graded latent activations may be ignored or misinterpreted by any downstream interpretability tool that assumes all features are independently interpretable and linear.
